\documentclass{article}
\usepackage{times}
\usepackage{a4wide}
\usepackage{latexsym}
\usepackage{mathrsfs}
\usepackage{amsmath}
\usepackage{amssymb}
\usepackage{amsthm}
\usepackage{ifthen}
\usepackage{stmaryrd}
\usepackage{algorithm}
\usepackage{ctex}
%\usepackage{algorithmic}
\usepackage{algpseudocode}
\usepackage{graphics}
\usepackage{epsfig}
\usepackage{setspace}
\setstretch{1}
%\usepackage{ramacros}

\addtolength{\textwidth}{20mm}
\addtolength{\oddsidemargin}{-10mm}
\addtolength{\textheight}{18mm}
\addtolength{\topmargin}{-9mm}


\newcounter{exercise}
\setcounter{exercise}{0}
\newcommand{\exercise}{
        \addtocounter{exercise}{1}
        \vspace{0.2in}
        \noindent
        {\bf \theexercise .}
        }

\newcommand{\REMARK}[1]{}


\newcommand{\NEWPART}{\vspace{.1in}
              \noindent}

\newcommand{\<}{
    \langle}

\renewcommand{\>}{
    \rangle}

\newcommand{\ceil}[1]{\left\lceil #1 \right\rceil}

\pagestyle{plain} \pagenumbering{arabic}

\title{{\bf Assignment 3} \\ {\large Deadline: Apr. 19}}

\author{125033910296 李博洋}
\date{}

\begin{document}
\maketitle

%{\large \noindent
%\begin{tabular}{lcl}
%{Jan 17 2007}.
%\end{tabular}
%}

{\large

\begin{exercise}
Prove the K\"onig theorem: Let $G$ be bipartite, then cardinality of maximum matching = cardinality of minimum vertex cover.
\end{exercise}
\\\\
\noindent\textbf{Solution:}
\newcommand{\card}[1]{\left|#1\right|} % 基数符号
\newcommand{\nuG}{\nu(G)} % 最大匹配基数
\newcommand{\tauG}{\tau(G)} % 最小顶点覆盖基数
\begin{proof}
设二分图 $G=(V,E)$，其顶点集 $V$ 可划分为两个互不相交的子集 $X$ 和 $Y$（二分划），满足所有边 $e\in E$ 都连接 $X$ 中的一个顶点和 $Y$ 中的一个顶点。

通过证明两个方向的不等式来完成定理的证明：
\[
\nuG \leq \tauG \quad \text{且} \quad \tauG \leq \nuG
\]

\paragraph{第一步：证明 $\nuG \leq \tauG$}
\\

设 $M$ 是 $G$ 的任意一个匹配，$C$ 是 $G$ 的任意一个顶点覆盖。
由于 $M$ 中的边两两无公共顶点，要覆盖这 $\card{M}$ 条边，$C$ 中至少需要 $\card{M}$ 个不同的顶点（每个顶点最多覆盖 $M$ 中的一条边）。

因此对所有匹配 $M$ 和顶点覆盖 $C$，都有 $\card{M} \leq \card{C}$。
取等号的极端情况可得：
\[
\nuG \leq \tauG
\]

\paragraph{第二步：证明 $\tauG \leq \nuG$}
\\

通过构造一个基数等于最大匹配基数的顶点覆盖来证明。设 $M$ 是二分图 $G$ 的一个最大匹配。
\begin{enumerate}
    \item 定义集合 $Z$：从 $X$ 中所有未被 $M$ 匹配的顶点出发，沿着交替路径（边交替属于 $M$ 和不属于 $M$ 的路径）能够到达的所有顶点的集合。
    \item 构造顶点覆盖：$C = (X \setminus Z) \cup (Y \cap Z)$。
\end{enumerate}

\subparagraph{验证 $C$ 是顶点覆盖}
假设存在一条边 $e=(u,v)$（$u\in X, v\in Y$），其两个端点都不在 $C$ 中，则：
\[
u \in X \cap Z \quad \text{且} \quad v \in Y \setminus Z
\]

我们分两种情况讨论边 $e$ 是否属于 $M$：
\begin{enumerate}
    \item 若 $e \notin M$：由于 $u \in Z$，说明存在从 $X$ 中未匹配顶点到 $u$ 的交替路径。在这条路径末尾添加边 $e$，就得到了到 $v$ 的交替路径，因此 $v \in Z$，与 $v \in Y \setminus Z$ 矛盾。
    \item 若 $e \in M$：由于 $u \in Z$，说明存在到 $u$ 的交替路径。因为 $e$ 是匹配边，到达 $u$ 的路径的最后一条边一定不是匹配边（否则路径会提前终止于 $v$）。在这条路径末尾添加匹配边 $e$，就得到了到 $v$ 的交替路径，因此 $v \in Z$，同样与 $v \in Y \setminus Z$ 矛盾。
\end{enumerate}

两种情况均矛盾，故不存在这样的边 $e$，即 $C$ 是 $G$ 的一个顶点覆盖。

\subparagraph{验证 $\card{C} = \card{M}$}
我们通过分析匹配边与集合 $Z$ 的关系来计算 $\card{C}$：
\begin{enumerate}
    \item $Y \cap Z$ 中的所有顶点都被 $M$ 匹配：若存在 $v \in Y \cap Z$ 未被匹配，则到达 $v$ 的交替路径是一条增广路径（起点是 $X$ 中未匹配顶点，终点是 $Y$ 中未匹配顶点），这与 $M$ 是最大匹配矛盾。
    \item $X \setminus Z$ 中的所有顶点都被 $M$ 匹配：若存在 $u \in X \setminus Z$ 未被匹配，则 $u$ 本身属于 $Z$（$Z$ 包含所有 $X$ 中未匹配顶点），与 $u \in X \setminus Z$ 矛盾。
    \item 匹配边在 $Z$ 内外的一一对应：
    \begin{itemize}
        \item 对于连接 $X \cap Z$ 和 $Y \cap Z$ 的匹配边：$X \cap Z$ 中被匹配的顶点与 $Y \cap Z$ 中被匹配的顶点一一对应，数量相等，记为 $k_1$。
        \item 对于连接 $X \setminus Z$ 和 $Y \setminus Z$ 的匹配边：$X \setminus Z$ 中所有顶点都被匹配，且其匹配点都在 $Y \setminus Z$ 中（否则匹配点在 $Y \cap Z$ 会导致该顶点属于 $X \cap Z$），数量为 $k_2 = \card{X \setminus Z}$。
    \end{itemize}
\end{enumerate}

因此，最大匹配的基数 $\card{M} = k_1 + k_2$。
而顶点覆盖的基数：
\[
\card{C} = \card{X \setminus Z} + \card{Y \cap Z} = k_2 + k_1
\]

故 $\card{C} = \card{M}$。这说明存在一个顶点覆盖 $C$，其基数等于最大匹配的基数，因此：
\[
\tauG \leq \nuG
\]

结合两个方向的不等式，我们得到：
\[
\nuG = \tauG
\]
\end{proof}

\begin{exercise}
Consider the following problem. You are given a flow network with unitcapacity
edges: It consists of a directed graph $G = (V, E)$, a source $s \in V$,
and a sink $t\in V;$ and $c_e = 1$ for every $e\in E$. You are also given a parameter $k$.
The goal is to delete $k$ edges so as to reduce the maximum $s$-$t$ flow in
$G$ by as much as possible. In other words, you should find a set of edges
$F \subseteq E$ so that $|F| = k$ and the maximum $s$-$t$ flow in $G' = (V, E-F)$ is as small
as possible subject to this.

Give a polynomial-time algorithm to solve this problem.
\end{exercise}
\\\\
\noindent\textbf{Solution:}
使用最大流算法（如Dinic算法）计算原图的最大s-t流$f$。
在最大流计算完成后，通过在残留网络中从s出发进行BFS，找到所有可达节点构成集合S，剩余节点构成集合T。
由于$G$是一个单位容量网络，因此从S指向T的所有边即为一个最小s-t割$C$，满足$|C| = f$。

此时，从最小s-t割$C$中删除$k$条边：
若$k \leq |C|$，选取任意$k$条边构成删除集$F$；
若$k > |C|$，选取$C$中所有$|C|$条边构成$F$（剩余$k - |C|$条边可任意选取，不影响结果）。

新图$G' = (V, E - F)$的最大s-t流即为\max(0, $|C| - k$)。

\begin{exercise}
 Find the minimum-cost flow subject to the constraint that the flow value is at
least 90\% of the maximum flow value.
\end{exercise}
\\\\
\noindent\textbf{Solution:}
\\
\noindent\textbf{步骤1：计算原图最大流量与流量下限}
\begin{enumerate}
    \item 使用任意最大流算法计算从$s$到$t$的最大流量$F_{\text{max}}$
    \item 计算流量下限：$F_{\text{min}}=0.9\times F_{\text{max}}$
    \item 计算允许的最大退流量：$\Delta_{\text{max}}=F_{\text{max}}-F_{\text{min}}$（最多可将流量从$F_{\text{max}}$减少多少）
\end{enumerate}

\noindent\textbf{步骤2：计算原图的最小费用最大流} \\
使用任意最小费用最大流算法（如连续最短路算法、原始-对偶算法）计算从$s$到$t$的最小费用最大流，得到：
\begin{itemize}
    \item 流量值$f=F_{\text{max}}$
    \item 总费用$C=C(F_{\text{max}})$
    \item 对应的残量网络$G_f$（包含正向边和反向边，反向边容量为已发送流量，单位费用为原边费用的相反数）
\end{itemize}

\noindent\textbf{步骤3：逐步退流优化总费用}\\
在残量网络中，每次选择\textbf{单位退流费用最小}的路径（即从汇点$t$到源点$s$的最短路径）退流，直到无法再降低总费用或达到流量下限。

初始化已退流量$\Delta=0$，循环执行：
\begin{enumerate}
    \item 在残量网络$G_f$中找到从$t$到$s$的最短路径$P$（单位费用之和最小）
    \item 若路径$P$的费用$\text{cost}(P)\geq0$：退流会增加总费用，此时总费用已达最小值，停止
    \item 计算路径$P$的最大可退流量$\Delta_p$：路径上所有边残量容量的最小值
    \item 计算本次实际退流量：$\Delta_{\text{actual}}=\min(\Delta_p, \Delta_{\text{max}}-\Delta)$（不能超过允许的最大退流量）
    \item 沿着路径$P$推送$\Delta_{\text{actual}}$个单位流量：
    \begin{itemize}
        \item 更新残量网络中各边的残量容量
        \item 总费用$C += \Delta_{\text{actual}}\times\text{cost}(P)$
        \item 总流量$f -= \Delta_{\text{actual}}$
        \item 已退流量$\Delta += \Delta_{\text{actual}}$
    \end{itemize}
    \item 若$\Delta\geq\Delta_{\text{max}}$：已达到流量下限$F_{\text{min}}$，不能再退流，停止
\end{enumerate}

\noindent\textbf{步骤4：输出结果}
最终得到的流$f$即为满足$|f|\geq 0.9\times F_{\text{max}}$的最小费用流，总费用为$C$。

\begin{exercise}
Given a directed network $G=(V,E)$ with the capacity $c(e)$ on each edge $e$. Associated with each node $v\in V$ is a demand $d_v$. If $d_v>0$, this indicates that the node $v$ has a demand of $d_v$ for flow; If $d_v<0$, this indicates that $v$ has a supply of $-d_v$. All capacities and demands are integers.
A circulation is a function $f$ that assigns a nonnegative real number to each edge and satisfies the following two conditions.
\begin{itemize}
\item[(i)]  (\textbf{Capacity conditions}) For each $e\in E$, we have $0\leq f(e)\leq c(e)$.
\item[(ii)]  (\textbf{Demand conditions}) For each $v\in V$, we have $f^{in}(v)-f^{out}(v)=d_v$.
\end{itemize}
Work out an algorithm to decide whether there exists a circulation that meets conditions (i) and (ii).
\end{exercise}
\\\\
\noindent\textbf{Solution:}
\\
\noindent\textbf{算法：转化为最大流问题}\\
首先判断，若$\sum_{v\in V} d_v \neq 0$，则不存在环流，否则将环流问题转化为源点-汇点的最大流问题，步骤如下：

\noindent\textbf{构造辅助图 \( G'=(V', E') \)}
\begin{enumerate}
    \item 新增源点 \( s \) 和汇点 \( t \)，令 \( V' = V \cup \{s, t\} \)；
    \item 保留原图所有边 \( e \in E \)，容量仍为 \( c(e) \)，即 \( E' \supseteq E \)；
    \item 对每个节点 \( v \in V \)：
        \begin{itemize}
            \item 若 \( d_v > 0 \)（需求节点）：添加边 \( (s, v) \) 到 \( E' \)，容量 \( c(s, v) = d_v \)；
            \item 若 \( d_v < 0 \)（供应节点）：添加边 \( (v, t) \) 到 \( E' \)，容量 \( c(v, t) = -d_v \)。
        \end{itemize}
\end{enumerate}

\noindent\textbf{计算最大流并判断}\\
定义总需求 \( D = \sum_{\substack{v \in V \\ d_v > 0}} d_v \)。在 \( G' \) 中计算从 \( s \) 到 \( t \) 的最大流 \( |f_{\text{max}}| \)：
\begin{enumerate}
    \item 若 \( |f_{\text{max}}| = D \)，则原图 \( G \) 存在满足条件的环流；
    \item 否则，不存在这样的环流。
\end{enumerate}

\begin{exercise}
Suppose that, in addition to edge capacities, a flow network has \textbf{vertex capacities}.
That is each vertex has a limit on how much flow can pass though. Show how to transform a flow network $G=(V,E)$ with vertex capacities into an equivalent flow network $G'=(V',E')$ without vertex capacities, such that a maximum flow in $G'$ has the same value as a maximum flow in $G$. How many vertices and edges does $G'$ have?
\end{exercise}

\begin{exercise}
SJTU is visited by prospective student groups. Let SJTU campus be modeled
as a weighted, directed graph $G$ containing locations $V$ connected by one-way
roads $E$. On a busy day, let $k$ be the number of campus tours that have to be
done at the same time. It is required that the paths of campus tours do not
use the same roads. Let the tour have $k$ starting locations $A = { a_1, a_2,..., a_k
} \in V$ . From the starting locations the groups are taken by a guide on a path
through $G$ to some ending location in $B =\{ b_1, b_2,..., b_k\} \in V$ . Your goal is
to find a path for each group $i$ from the starting location, $a_i$
, to any ending
location $b_j$ such that no two paths share any edges, and no two groups end in
the same location $b_j$ .


\begin{itemize}
\item Design an algorithm to find $k$ paths $a_i \rightarrow^* b_j$ that start and end at different
vertices and such that they do not share any edges.


\item Modify your algorithm to find $k$ paths $a_i \rightarrow^* b_j$ , that start and end in different
locations and such that they share neither vertices nor edges.

\end{itemize}


\end{exercise}


}

\end{document}
